\section{Guénon lecteur de Nietzsche}

\begin{quotex}
Mais seule la littérature peut vous donner cette sensation de contact avec un autre esprit humain, avec l'intégralité de cet esprit, ses faiblesses et ses grandeurs, ses limitations, ses petitesses, ses idées fixes, ses croyances ; avec tout ce qui l'émeut, l'intéresse, l'excite ou lui répugne. Seule la littérature peut vous permettre d'entrer en contact avec l'esprit d'un mort, de manière plus directe, plus complète et plus profonde que ne le ferait même la conversation avec un ami --- aussi profonde, aussi durable que soit une amitié, jamais on ne se livre, dans une conversation, aussi complètement qu'on ne le fait devant une feuille vide, s'adressant à un destinataire inconnu.
\end{quotex}


Dans son roman, \emph{Soumission}, \textbf{Michel Houellebecq} se réfère à une thèse de philosophie, soutenue à l'université catholique de Louvain-la-Neuve, signée Robert Rediger, et intitulée \emph{Guénon lecteur de Nietzsche}. Au cours des derniers mois, j'ai fait de nombreux efforts pour traquer une copie. Une amie, employée du Département d'État en poste à Bruxelles, a réussi à localiser une photocopie pour moi. Quand elle était en route pour visiter sa famille aux États-Unis pour Noël, j'ai pu l'inspecter pendant son escale à l'aéroport. Bien que je n'ai pas eu beaucoup de temps, je cherchai à en extraire les principaux thèmes.

La première chose à souligner est l'influence profonde que la thèse de Rediger a eu sur Houellebecq lui-même. Ceci est représenté dans les principaux thèmes du roman tel qu'ils sont exprimés par la voix de François, son \emph{alter ego}. En outre, le mémoire fait ressortir plusieurs thèmes rudimentaires qui ont été complétés dans les écrits ultérieurs de Rediger. De toute évidence, René Guénon n'a jamais été un «  lecteur » de Nietzsche. Rediger est arrivé à cette conclusion quand il a écrit:

\begin{quotex}
Guénon, à bien y réfléchir, n'a pas été tant que ça influencé par Nietzsche ; son rejet du monde moderne est tout aussi fort, mais il vient de sources radicalement différentes.
\end{quotex}


\paragraph{La Décadence}


Ce ne sont donc pas les sources qui sont d'intérêt, mais plutôt comment Guénon et Nietzsche sont parvenus à des conclusions similaires. Plus précisément, il en est la reconnaissance de la décadence de l'Europe, l'échec du christianisme à la prévenir, et la montée de l'Islam. François note ceci:

\begin{quotex}
Il fallait se rendre à l'évidence : parvenue à un degré de décomposition répugnant, l'Europe occidentale n'était plus en état de se sauver elle-même -- pas davantage que ne l'avait été la Rome antique au cinquième siècle de notre ère. L'arrivée massive de populations immigrées empreintes d'une culture traditionnelle encore marquée par les hiérarchies naturelles, la soumission de la femme et le respect dû aux anciens constituait une chance historique pour le réarmement moral et familial de l'Europe, ouvrait la perspective d'un nouvel âge d'or pour le vieux continent.
\end{quotex}


Les États-Unis sont également dans un état de décomposition, même si les circonstances particulières sont très différentes. La principale dif\-fé\-rence, bien sûr, est que l'européen reconnaît plus ou moins sa situation alors que son homologue américain est allègrement inconscient de sa situation réelle. Même ceux qui se disent conservateurs aux Etats-Unis vivent dans l'obscurité. Les talk-shows d'extrême droite sont contrariés par la violence musulmane mais ils, de toutes les personnes, devraient être plus préoccupés par ceux qui ne peuvent tuer l'âme plutôt que ceux qui tuent le corps. Par conséquent, ils condamnent l'islam pour tout dans ce qui est traditionnel. Par exemple, ils soulignent la préoccupation de l'Occident pour les femmes et les droits des homosexuels, tout en louant la séparation de toute autorité spirituelle sur les questions politiques.

Rediger a ironiquement souligné:

\begin{quotex}
Il était tragique, plaidait-il avec ferveur, qu'une hostilité irraisonnée à l'islam les empêche de reconnaître cette évidence : ils étaient, sur l'essentiel, en parfait accord avec les musulmans. Sur le rejet de l'athéisme et de l'humanisme, sur la nécessaire soumission de la femme, sur le retour au patriarcat : leur combat, à tous points de vue, était exactement le même.
\end{quotex}


\paragraph{Le Patriarcat}


Cela se voit tout d'abord dans la fin du patriarcat dans l'Ouest. Dans les mots de François, presque une citation directe du mémoire:

\begin{quotex}
Le patriarcat avait le mérite minimum d'exister, enfin je veux dire en tant que système social il persévérait dans son être, il y avait des familles avec des enfants, qui reproduisaient en gros le même schéma, bref ça tournait ; là il n'y a plus assez d'enfants, donc c'est plié.
\end{quotex}


Le libéralisme a pu renverser toutes les institutions établies: les églises, l'éducation, le travail, dirigeants politiques, et ainsi de suite. Cependant, l'attaque de la famille ne peut pas être maintenue, car sans les enfants, aucune société ne peut continuer à exister.

Néanmoins, le libéralisme poursuit catégorique, alors que le taux de natalité chutent. Il n'y a alors plus une communauté organique mais plutôt un tas de consommateurs indépendants. François énumère quelques exemples:

\begin{itemize}


\item le bobo éco-responsable

\item la bourgeoise show-off

\item la clubbeuse gay-friendly

\item le satanic geek

\item le techno-zen
\end{itemize}


En d'autres termes, tout le monde est un « type », en essayant trop dur de se doter d'une identité sui generis. Le lecteur en peut identifier plus. Par exemple, il y a un groupe Facebook appelé « libéral et fier de l'être ». Ainsi, il n'est plus question d'une vision bien raisonnée du monde, mais plutôt un marqueur de l'identité tribale.

François, malgré lui, est un ``type'', peut-être le professeur vieillissement ou l'intellectuel décadent. Sa propre vie illustre la stérilité du monde moderne. Rejetant la vie de famille, il s'accouple avec ses élèves, ou avec des call-girls lorsque cela est nécessaire. Les scènes de sexe ne sont pas destinées à être érotique, mais plutôt triste. Sa préférence est pour le sexe anal ou bien il aime éjaculé dans la bouche d'une femme -- n'importe où sauf où un homme bien élevé envisagerait comme saine et normale.

\paragraph{L'assimilation et la transformation}


Rediger prend alors un thème mentionné il y a quelques années sur Gornahoor dans The Prophecies of Guénon\footnote{\url{https://www.gornahoor.net/?p=35}}. Dans \emph{La crise du monde moderne}, Guénon écrit que seul le catholicisme pourrait restaurer la Tradition monde moderne. Rediger est d'accord avec Guénon sur la tradition médiévale:

\begin{quotex}
La chrétienté médiévale avait été une grande civilisation, dont les accomplissements artistiques resteraient éternellement vivants dans la mémoire des hommes ; mais peu à peu elle avait perdu du terrain, elle avait dû composer avec le rationalisme, renoncer à se soumettre le pouvoir temporel, ainsi peu à peu elle s'était condamnée.
\end{quotex}


Selon Guénon, les dirigeants de l'Eglise auraient oublié les vérités les plus profondes de la tradition. Rediger va plus loin, affirmant que l'Eglise elle-même est activement engagée dans la subversion de la Tradition. Il écrit:

\begin{quotex}
À force de minauderies, de chatteries et de pelotage honteux des progressistes, l'Église catholique était devenue incapable de s'opposer à la décadence des moeurs. De rejeter nettement, vigoureusement, le mariage homosexuel, le droit à l'avortement et le travail des femmes.
\end{quotex}


Même Guénon, après avoir échoué d'avoir un impact sur la pensée catholique, ouvertement converti à l'Islam. Bien que Guénon a affirmé qu'il l'a fait pour des raisons personnelles, Rediger voit une perspective plus large:

\begin{quotex}
Ce combat nécessaire pour l'instauration d'une nouvelle phase organique de civilisation ne pouvait plus, aujourd'hui, être mené au nom du christianisme ; c'était l'islam, religion soeur, plus récente, plus simple et plus vraie (car pourquoi Guénon par exemple s'était-il converti à l'islam ? Guénon était avant tout un esprit scientifique, et il avait choisi l'islam en scientifique, par économie de concepts ; et pour éviter, aussi, certaines croyances irrationnelles marginales, telles que la présence réelle dans l'Eucharistie),
\end{quotex}


Comme l'avait prédit Guénon, les choix possibles sont la dégénérescence, l'assimilation et la transformation. La dégénérescence ne peut pas continuer sans relâche. Dans le roman, Houellebecq opte pour une assimilation douce et puis une transformation. Il prévoit une alliance entre la gauche et les éléments islamistes. Cela ne veut pas si farfelu puisque le libéralisme croit qu'il peut utiliser l'islam à éroder l'influence chrétienne. Toutefois, dans le scénario de Houellebecq, il est l'islam qui assimile la gauche.

\paragraph{La polygamie et l'ordre social}


Le seul point d'importance pour Guénon portait sur les éléments métaphysiques; tellement qu'il n'offre que peu ou pas d'aperçu de la vie religieuse musulmane réelle et pratique. Rediger, d'autre part, développe sur ces influences tandis que, dans un même temps, il y ajoute les éléments distinctifs de l'Ouest à l'Islam. Pour cette partie de la thèse, il repose davantage sur Nietzsche que sur Guénon.

Puis, Rediger, également suivant Nietzsche, affirme que le christianisme est fondamentalement une religion féminine. Rediger cite Nietzsche de L'antéchrist:

\begin{quotex}
Si l'islam méprise le christianisme, il a mille raisons pour cela ; l'islam a des hommes pour condition première.
\end{quotex}


Rediger attribue les idéaux de l'humanisme et les droits de l'homme au dogme de l'Incarnation. Maintenant, vous pourriez objecter que le christianisme médiéval, qui Rediger a loué, était lui-même patriarcal et masculin. Quelqu'un peut-il accuser Constantin, Clovis, Charlemagne, Arthur, ou Boucicaut\footnote{\url{https://www.gornahoor.net/?p=277}} d'être féminin ? En outre, les médiévaux ont admiré les grands guerriers païens comme Hector, Alexandre, Scipion, et Jules César. Même St Paul a considéré mollesse comme un péché. À ce stade, Rediger fait la seule référence à Julius Evola que j'ai remarqué: il affirme aussi que cet élément masculin dans la chrétienté médiévale était due uniquement à des vestiges encore existants du paganisme, mais pas quelque chose de spécifiquement chrétien.

Tout d'abord, Rediger accepte la classification de Nietzsche d'affirmation de la vie et religions de la vie de renoncement, dont l'Islam appartient à la première. Il écrit:

\begin{quotex}
l'islam accepte le monde, et il l'accepte dans son intégralité, il accepte le monde tel quel, pour parler comme Nietzsche. Le point de vue du bouddhisme est que le monde est dukkha -- inadéquation, souffrance. Le christianisme lui-même manifeste de sérieuses réserves -- Satan n'est-il pas qualifié de ``prince de ce monde'' ? Pour l'islam au contraire la création divine est parfaite, c'est un chef d'œuvre absolu. Qu'est-ce que le Coran au fond, sinon un immense poème mystique de louange ? De louange au Créateur, et de soumission à ses lois.
\end{quotex}


Rediger se tourne de nouveau à Nietzsche pour justifier l'ordre hié\-rar\-chique et patriarcal islamique. Il y aurait un petit groupe d'aristocrates qui surplomberaient une large base de gens ordinaires. Destitution serait réduite en raison de l'exigence de l'aumône. En outre, la famille élargie serait les « premiers intervenants », pour ainsi dire, dans le cas des tragédies de la vie. Ces éléments réduiraient la taille du budget de l'aide sociale. De plus, le chômage serait réduit parce que les femmes auraient quittés le milieu du travail pour rester à la maison, pour l'éducation des enfants.

La polygamie, aussi, est comprise comme la volonté de puissance. Le plus fort et le plus génétiquement ajustement aurait plus de femmes, d'où plus d'enfants. Bien sûr, cela signifie que certains hommes n'auraient pas de descendance du tout. Cela me rappelle ce que nous a récemment écrit à propos de Social Surgery\footnote{\url{https://www.gornahoor.net/?p=8253}}. Dans quelle mesure la communauté a le droit de réglementer la composition génétique de ses membres, en supposant que l'objectif soit le bien commun ?

À ce stade, Rediger devient moins convaincant car il ne pousse pas ses idées jusqu'à leur conclusion logique. Il mentionne avec raison que les hommes de tous types et castes choisiront le même genre de femmes. Par exemple, dans le choix du partenaire, la plupart des hommes, in\-dépen\-dam\-ment de l'intelligence ou de la station dans la vie, préféreraient une Taylor Swift. Les femmes, d'autre part, sont plus malléables dans leur choix du partenaire; par exemple, elles peuvent négliger l'apparence ou l'âge d'un homme s'il est riche ou puissant.

En outre, l'attraction de la femme peut être formée. Rediger, donc, veut définir les intellectuels comme des ``mâles dominants'', qui méritent des femmes de meilleure qualité. Il serait incroyable que, dans l'Ouest comme il en est maintenant, les femmes seront attirés par les thomistes studieux plutôt que le «  star  » de l'équipe. Cependant, en y réfléchissant, les étudiantes sont souvent attirées par leurs homologues masculins. Les femmes asiatiques, et dans une moindre mesure les Latines, sont plus susceptibles que les femmes européennes de trouver un grand ingénieur sexuellement désirable.

Pour ce faire, Rediger affirme qu'il y aurait des femmes qui serviraient d'entremetteuses. Les détails de ce processus sont flous. Il est facile de déterminer si les pays islamiques sont biologiquement plus en forme que les pays occidentaux. Ma compréhension est que le contraire est parfois le cas en raison de la préférence pour les mariages entre cousins. Mais peut-être, la version occidentalisée nietzschéenne de l'Islam de Rediger deviendrait la norme.

\paragraph{Avec un gémissement}


Dans l'ensemble, la structure sociale est basée sur la soumission, à commencer par la soumission à Dieu. Ensuite, il y a la soumission à l'ordre social, et enfin la soumission des femmes aux hommes. Houellebecq voit dans \emph{l'Histoire d'O} de Pauline Réage le modèle de soumission. Il voit cela comme se produisant sans grand incident. Ceci est difficile à croire, mais, étant donné que les « 50 Shades of Grey » a vendu plus de 70 millions d'exemplaires, peut-être il y a un désir secret d'être dominée.

Ainsi, la transition se termine sans éclat, mais avec un gémissement. Pourquoi en serait-il de cette façon? Certains diront que ce n'est qu'une question de commodité, que les intellectuels et les autres profitent des nouvelles possibilités qui s'ouvrent pour eux. Rediger se termine par une réponse différente:

\begin{quotex}
Au fond, c'était un mystère ; Dieu en avait décidé ainsi.
\end{quotex}


Si tel est le cas, alors le mystère des trois anneaux\footnote{\url{https://www.gornahoor.net/?p=7872}} aura été résolu.\footnote{For an English translation, see Section \ref{sec:GuenonReader} in this book.}

H/T DB for proofreading the French text.


\begin{center}* * *\end{center}

\begin{footnotesize}
\begin{sffamily}

\texttt{Mark Citadel on 2016-02-17 at 10:38 said:}

Rediger seems to echo Evola, but with a much more pro-Islamic tone (this makes sense as in the book, he has benefited from Islamic rule). I continue to believe Guenon was right on the viability for Christianity, though certainly the virility of Catholicism has slipped since his day. Understandably, he never interacted with Orthodoxy as most Western scholars would have presumed it dead, especially after the Iron Curtain descended.

\end{sffamily}
\end{footnotesize}

